\subsection{Receiving and Decoding the Pose Stream}\label{subsec:receiver}

\texttt{ManoLiveReceiver} is the counterpart of \texttt{mano\_unity\_streamer}
(\autoref{subsec:streaming}) and reverses each of its steps. On start it opens a
receiver on the same port \texttt{5052} and runs the receive loop on a dedicated
background thread. Each arriving datagram is decoded from UTF-8 and stored as the
latest message under a lock, so the socket never blocks the main thread. This
mirrors the decoupling used on the capture side~(\autoref{subsec:hand_tracking}):
just as the Python tracker reads whatever detection is currently available rather
than waiting on the detector, the Unity main loop acts on whatever datagram has
most recently arrived. Because each pose supersedes the last, only the newest
message is retained; a datagram that arrives while the previous one is still
unread is simply overwritten, which is the intended behaviour for the
fire-and-forget stream of~\autoref{subsec:udp}.

Once per frame, \texttt{Update} takes the buffered message under the lock, clears
it, and deserialises it with \texttt{Newtonsoft.Json} into a payload whose fields
mirror the datagram schema of~\autoref{fig:udp-schema}.
Any field whose hand was absent in that frame is null and is skipped.

The rotations are decoded by \texttt{ApplyPose}, which inverts the axis-angle
encoding produced upstream~(\autoref{subsec:mesh_recon}). Each of the sixteen
rotations is a vector whose direction is the axis of rotation and whose magnitude
is the angle in radians; the method recovers the angle as the vector's length and
the axis as its normalisation, and reconstructs the rotation through Unity's
\texttt{Quaternion.AngleAxis}. The sixteen quaternions are written as local
rotations onto a sixteen-joint armature ordered as the wrist followed by the three
joints of each finger --- the \acrshort{mano} joint convention
of~\autoref{subsec:anatomy}. The first rotation is the global orientation of the
wrist; the remaining fifteen are the per-finger joint rotations, exactly the
decomposition the solver emitted into \acrshort{mano} pose-vector
order~(\autoref{subsec:ik_impl}).

Placement is handled separately by \texttt{ApplyAnchor}. The world-landmark
skeleton is wrist-relative and so carries no information about camera distance,
which is why the streamer supplies a separate anchor combining the wrist's
image-space position with a monocular depth estimate~(\autoref{subsec:landmark_smoothing}).
The receiver scales this anchor into scene units and inverts one axis to match
Unity's left-handed coordinate frame, then writes it to the position of the
armature root. Rotation and position are therefore restored independently: the
joint rotations pose the fingers, while the anchor places the whole hand where it
sits relative to the camera.

One mapping inside the receiver is deliberately crossed. The capture loop mirrors
each frame horizontally to present a selfie-style view~(\autoref{subsec:streaming}),
which leaves the detector's left and right handedness labels reflected with respect
to the user's actual hands. The receiver compensates by routing the payload's
right-hand pose, gesture, and anchor to its left-hand avatar, and the left-hand
fields to its right-hand avatar.

The receiver is also the point at which software latency is measured. Each
frame, it subtracts the datagram's send timestamp from the current time to obtain
the software latency of that frame, and, when logging is enabled, appends the
value to a CSV file over a fixed window. This is the instrument behind the headline
timing result of this thesis~(\autoref{subsec:contributions}). For the present
chapter the receiver's role is complete once it has produced, every frame, two
posed and placed hand skeletons together with their gesture labels.

Beyond decoding, the receiver exposes precisely what every scene consumes: the two
bone arrays, the two armature roots, and the current gesture for each hand. It also
holds references to the two hand meshes, so that they can be hidden when the hands
are driving an avatar rather than appearing as themselves. No interaction logic
lives here; that is the subject of the remaining sections.

\begin{figure}[H]
	\centering
	\captionsetup{justification=centering}
	\includegraphics[width=0.9\linewidth, frame]{img/unity_live_hands.png}
	\caption{The two \acrshort{mano} hand meshes in Unity,
	         posed and placed live from the decoded stream.}
	\label{fig:unity-hands}
\end{figure}
